\documentclass[11pt,letterpaper]{article}
\usepackage[margin=1in]{geometry}\usepackage{amsmath,amssymb,booktabs,array,microtype}\usepackage[T1]{fontenc}\usepackage{lmodern}
\usepackage[colorlinks=true,linkcolor=blue,citecolor=blue,urlcolor=blue]{hyperref}
\setlength{\parindent}{0pt}\setlength{\parskip}{5pt}\sloppy
\title{A Photocount Reading of the Interpolating Function,\\ and a No-Go for Deriving Its Normalisation from a Clock}
\author{Carl P.\ Zimmerman\\ \small Briar Creek Tech\\ \small AI-assisted research programme; not peer reviewed}
\date{13 September 2026}
\begin{document}\maketitle
\begin{abstract}
A companion paper removed the MOND acceleration scale as an independent parameter and found that 155 SPARC rotation curves then select, with nothing fitted, the member $n=2$ of the family $\mu_n(Y)=1-(1+Y)^{-n}$, which is $\kappa=1/2$. This paper reports what happened when we tried to derive that integer, and one structural result obtained along the way. \textbf{First, a no-go.} A cuscuton clock's own field equation fixes $V'(\tau)=-3HU$, so a vanishing potential forces a vanishing expansion rate: any cuscuton clock in an expanding universe \emph{requires} a potential. A potential is precisely the free additive constant that makes the acceleration-scale normalisation a zero mode. Hence the twelve-gate relativistic construction of our earlier work provably cannot derive $\kappa$, and the two tracks cannot be merged --- they assign values fourteen orders apart to the same coefficient, and forcing agreement demands a margin exceeding one. \textbf{Second, an identification.} The family is exactly Mandel's $n$-mode photocount formula: for $n$ independent thermal modes of mean occupancy $Y$, $\mu_n$ is the probability that at least one quantum is present. Equivalently it is one minus a Tsallis $q$-exponential with $q=1+1/n$, the ordinary exponential kernel being the infinite-mode limit. This explains why the exponent is an \emph{integer} --- it is a mode count --- without fixing its value, and replaces the question with a sharper one: why does the vacuum present exactly two modes to an acceleration? \textbf{Third, the answer to that is not available, and we report the closures.} The modes are not Unruh modes: a fixed-frequency thermal occupancy is exponentially small where the reading requires it linear, failing by 41 orders at $Y=10^{-2}$. Twelve pre-registered mathematical criteria scatter over three different members, with $n=1$ the modal choice; only the convergence boundary selects $n=2$, and as a boundary rather than an optimum. Effective-field-theory positivity is inapplicable (the kinetic matrix degenerates where $\mu(0)=0$) and causality is saturated identically by every kernel. No published entropic or holographic derivation produces the family beyond $n=1$. Finally, the $d$-dimensional generalisation $\mu_n^{(d)}=1-(1+Y^{d-2})^{-n}$ exists and is forced, but the mode count enters the deep limit as a multiplicative prefactor and cannot change a power --- so it is dimensionally inert and no dimension fixes it. \textbf{We conclude that the exponent is an empirical quantity, and that this is the result rather than a placeholder for one.}
\end{abstract}

\section{A no-go for the clock}
Our earlier relativistic construction \cite{paper26} carries a cuscuton clock $\tau$ with action $\int\!\sqrt{-g}\,[\,s\,W(Y,\tau)-V(\tau)+\dots]$, $s=\sqrt{-\nabla\tau\cdot\nabla\tau}$. Varying $\tau$ at constant $U\equiv W(0,\tau)$ gives $U\,\nabla_\mu n^\mu=-V'(\tau)$, and in a Friedmann background $\nabla_\mu n^\mu=3H$, so
\begin{equation}
V'(\tau)=-3HU .
\label{clock}
\end{equation}
A vanishing potential therefore forces $H=0$: \textbf{any cuscuton clock in an expanding universe requires a potential.}

This matters because of a zero mode. For any \emph{free} interpolating function $F$ of derivative invariants, the shift $F\to F+c$ leaves $F'$ untouched --- and $F'$ is all the scalar equation and the static force law ever see --- while shifting the metric-multiplying piece by exactly $c$, which is a cosmological constant and nothing else. The normalisation relating $a_0$ to the vacuum energy is therefore unfixable whenever the action contains a free additive constant, \emph{regardless} of field content \cite{paper27}. Equation (\ref{clock}) says the clock construction cannot dispense with one. It is therefore inside the class the no-go covers, and no further work on it will produce $\kappa$.

The two tracks also cannot be merged directly. Both call $W(0)$ the same coefficient $U$; the clock construction gives $U=m_{\rm rel}\rho_{\rm dm}=3.68\times10^{-25}\,$eV$^4$ at its parameter point, while the parameter-free reading gives $U=\rho_\Lambda=2.52\times10^{-11}\,$eV$^4$, a ratio of $6.9\times10^{13}$. Forcing agreement needs $m_{\rm rel}=\rho_\Lambda/\rho_{\rm dm}=2.58$, and $m_{\rm rel}=1-\mu$ cannot exceed unity. The merge is arithmetically blocked, not merely strained. \emph{The one loophole:} (\ref{clock}) is written at constant $U$; with $U$ depending on $\tau$ there is an extra term, and whether it can substitute for the potential is not computed here.

An architecture that is both relativistic and $\kappa$-deriving must therefore have its expansion driven by the \emph{same} function that carries the gradient sector, with no independent potential. That rules out the cuscuton as the timekeeper.

\section{The family is a photocount formula}
A single bosonic mode in thermal equilibrium with mean occupancy $Y$ has the geometric distribution $P(k)=Y^k/(1+Y)^{k+1}$, so it is empty with probability $1/(1+Y)$. For $n$ independent such modes,
\begin{equation}
\mu_n(Y)=1-(1+Y)^{-n}=P(\text{at least one quantum present}),
\end{equation}
which is Mandel's $n$-mode photocount formula. We verified the identification symbolically and by Monte Carlo (400{,}000 trials at $Y=2$, $n=2$: $0.88914$ against the closed form $0.88889$). Under the principle of \cite{paper27}, $Y$ is the acceleration in units of the dark-energy acceleration, so the reading is direct: \emph{the response to an acceleration is the probability that it excites at least one quantum out of the vacuum.}

The same family is one minus a Tsallis $q$-exponential, $\mu_n=1-\exp_q(-nY)$ with $q=1+1/n$, so the mode count and the deformation parameter are one quantity, and the exponential kernel in common use is the infinite-mode limit --- the case in which the mode count has been forgotten. The composition law $1-\mu_n=(1-\mu_1)^n$ becomes obvious: $n$ independent modes.

This buys the \emph{form} and not the \emph{value}. It explains why the exponent is an integer, and replaces the question with: \textbf{why does the vacuum present exactly two independent modes to an acceleration?}

\section{The modes are not Unruh modes}
The obvious answer is that an accelerating body sees a bath at $T=g/2\pi$ and the modes are that bath's. A mode of frequency $\omega=s/2\pi$ does give occupancy $g/s$ --- but only in the classical, high-occupancy limit. The photocount reading needs occupancy \emph{exactly} $Y=g/s$ at all accelerations, and a thermal occupancy is exponentially small at low temperature:

\begin{center}
\begin{tabular}{ccc}
\toprule
$Y=g/s$ & reading requires & thermal gives \\
\midrule
$1$ & $1$ & $0.582$ \\
$10^{-1}$ & $10^{-1}$ & $4.5\times10^{-5}$ \\
$10^{-2}$ & $10^{-2}$ & $3.7\times10^{-44}$ \\
\bottomrule
\end{tabular}
\end{center}

Forty-one orders, in the only regime MOND is about. The modes are not Unruh modes, nor any fixed-frequency thermal modes. This matches independently the fact that every thermal-screen derivation in the literature has \emph{exponentially small} deep-MOND corrections where this family's are analytic. What survives is that the occupancy is exactly linear in the acceleration --- an equipartition statement, how many quanta of the dark-energy acceleration fit into this one, not a Boltzmann one.

\section{Four closed routes}
\emph{Pre-registered mathematical criteria.} Twelve natural distinguishing properties of the family were fixed before evaluation, each statable without reference to the number two. Nine gave a unique selection, scattering over three members: $n=1$ takes five, $n=3/2$ two, $n=2$ two. The modal choice is $n=1$. The two criteria selecting $n=2$ --- convergence of $\int(1-\mu_n)$ and finiteness of $\lim(F-z)$ --- are the same fact in two dressings, and select it as a \emph{boundary} rather than an interior optimum. With nine criteria over three members, chance places three on any one. The one substantive residue is that $n=2$ is the \emph{marginal} index, which is what a criticality argument would want; that is recorded as suggestive and not more.

\emph{Effective field theory.} Positivity bounds are inapplicable: with $\mu(0)=0$ the fluctuation kinetic matrix degenerates at zero gradient, so there is no propagator and no $S$-matrix, and AQUAL is not a derivative expansion. Granting them anyway, the single scale-free invariant is $9(n+1)/(16n)$, continuous and monotone, so a two-sided bound could exclude a range but never select an integer. Causality is non-discriminating: around a static gradient $c_\parallel^2=1+d\ln\mu/d\ln Y$, which tends to $2$ in the deep regime for \emph{every} member of the family, and indeed for any kernel with a deep-MOND limit, since $\mu\sim Y^p$ gives $c_\parallel^2\to p+1$ and MOND forces $p=1$. Nobody has aimed positivity or causality at MOND interpolating functions before.

\emph{Existing derivations.} No published entropic or holographic derivation produces this family beyond the $n=1$ member, which modified-R\'enyi entropic gravity gives exactly. Every factor of two in that literature sits in the acceleration scale, never in the exponent.

\emph{Dimension.} Keeping the photocount structure, the dimension can only enter through the occupancy, giving $\mu_n^{(d)}(Y)=1-(1+Y^{d-2})^{-n}$, which reduces to the family at $d=3$ and carries Milgrom's conformally invariant deep power in general. But in the deep limit the response is count $\times$ occupancy, so the count is a multiplicative prefactor and cannot change a power: $d\ln(\text{response})/d\ln Y=d-2$, free of $n$. The count is \emph{dimensionally inert}, so no dimension fixes it. Four candidate counts --- transverse directions, graviton polarisations, screen dimensions, and the two branches of the gradient invariant --- all give two at $d=3$ and the generalisation separates none of them. A by-product: the equipartition reading of the occupancy is specifically a three-dimensional fact.

\section{Conclusion}
The exponent is an empirical quantity. The rotation curves prefer $n=2$ over $n=1$ by $0.016$ dex \cite{paper27}; sharpening that means beating the relation's $\sim0.045$ dex intrinsic scatter by a factor of a few, which needs better mass-to-light ratios and distances rather than better theory. After the closures above we do not regard this as a placeholder for a derivation that will appear shortly. \textbf{It is the result.}

\section{What is not claimed}
The clock no-go is written at constant $U$ and the $\tau$-dependent case is not computed. The photocount identification is an \emph{identification}: nothing here shows the vacuum response \emph{is} such a count, only that the function the data selected has that form; and the analytic-versus-exponential mismatch cuts against reading it as a literal thermal screen. Twelve criteria is not exhaustive. The $d$-dimensional deep power is Milgrom's conformal argument, quoted. Uniqueness of the generalisation holds within the photocount reading only. The Unruh test uses a single fixed frequency; a distribution of frequencies could give linear occupancy, but a Gamma-distributed rate is exactly what generates this family, so that route risks assuming the answer. This paper concerns one coefficient and is not a verdict on any construction.

\section{Reproducibility}
\texttt{fable\_independent\_2026/} lanes \texttt{L235}--\texttt{L239} with committed outputs and results files, and \texttt{hunt\_2026/eft01\_positivity\_causality\_mu\_family\_2026.py}. Lean file \texttt{lean\_2026/Mondlean.lean}: 160 theorems, exit 0, zero \texttt{sorry}.

\begin{thebibliography}{9}
\bibitem{paper26} C.~P.~Zimmerman, \emph{A Conserved Clock Rate, and the Board Run End to End}, doi:10.5281/zenodo.22731066 (2026).
\bibitem{paper27} C.~P.~Zimmerman, \emph{A Parameter-Free Radial Acceleration Relation from the Dark-Energy Scale}, doi:10.5281/zenodo.22731370 (2026).
\bibitem{fm12} B.~Famaey, S.~McGaugh, Living Rev.\ Relativity \textbf{15}, 10 (2012), arXiv:1112.3960.
\bibitem{mandel} L.~Mandel, E.~Wolf, \emph{Optical Coherence and Quantum Optics}, Cambridge (1995).
\bibitem{cuscuton} N.~Afshordi, D.~J.~H.~Chung, G.~Geshnizjani, Phys.\ Rev.\ D \textbf{75}, 083513 (2007).
\bibitem{bruneton} J.-P.~Bruneton, G.~Esposito-Far\`ese, Phys.\ Rev.\ D \textbf{76}, 124012 (2007), arXiv:0705.4043.
\bibitem{repo} The research repository, \url{https://github.com/carlzimmerman/zimmerman-formula} (2026).
\end{thebibliography}
\end{document}
